\documentclass[11pt]{article}
\usepackage[margin=1in]{geometry}
\usepackage{amsmath,amssymb,booktabs,hyperref,siunitx}
\usepackage{xcolor}
\newcommand{\pass}{\textcolor{green!55!black}{\textbf{PASS}}}
\newcommand{\partialv}{\textcolor{orange!85!black}{\textbf{PARTIAL}}}
\newcommand{\oos}{\textcolor{gray}{\textbf{out of scope}}}
\title{Replication Report: \\ \emph{Current Driven Magnetization Dynamics in Helical Spin Density Waves}\\ \normalsize arXiv:cond-mat/0511224 (Wessely, Skubic, Nordstr\"om; PRB \textbf{73}, 144431 (2006))}
\author{REPLICATE-PROJECT / TEXTURES-100 automated replication}
\date{\today}
\begin{document}
\maketitle

\section*{Classification audit (up front)}
This paper was filed under the \textbf{loop-current} texture class. It is
\textbf{not} a loop-current / kagome flux-phase paper: there is no orbital
current, no Peierls flux, no kagome lattice, and no time-reversal-breaking
kinetic term. The physics is \textbf{spin-transfer torque (STT) driven
magnetization dynamics in a helical spin density wave (spin spiral)}: a charge
current along the spiral axis transfers spin to the local moments and rigidly
rotates ("slides") the spiral. The shared kagome loop-current kernel
(\texttt{loop\_current\_kagome\_kernel.py}) is therefore \textbf{not applicable}
and was \textbf{not imported}. We replicate the actual in-scope core with a
purpose-built spin-spiral tight-binding model. \emph{This is a taxonomy
mislabel in the input set, reported honestly, not a physics failure.}

\section{Paper summary}
The authors propose that a DC charge current flowing along the axis of a
helical spin density wave (SDW) exerts a bulk spin-transfer torque that rigidly
rotates the spiral. Using semiclassical Boltzmann linear response (their
Eqs.~3--6) combined with a full-potential APW+lo LSDA calculation of Er
($q=0.20\,(2\pi/c)$, moment $\parallel[100]$), they obtain a torque--current
tensor
\begin{equation}
  C=\hbar\begin{pmatrix}0&0&0\\0&0&0.5\\0&0&0\end{pmatrix}\;\si{\angstrom^2},
\end{equation}
i.e.\ only a current \emph{along} the spiral axis produces a
rotate-the-spiral torque. A current density of \SI{e7}{A/cm^2} yields a spiral
precession of \SI{0.07}{GHz}; a crude adiabatic estimate gives $\sim 4\times$
this value. They further note a Fermi-surface parallel polarization
$P\approx-0.5$ (conduction spins tilted \ang{30} from the axis) and propose an
oscillating-conductance nanopillar readout.

\section{What we replicate}
The absolute Er numbers derive from DFT we do not run (offline / free-endpoints
scope). We instead verify the \emph{convention-independent} claims with a
minimal first-principles-free realization of the same mechanism: a 1D
spin-spiral $s$-band along the spiral axis, spin-$\tfrac12$, exchange field
$\Delta$ rotating in-plane. Via the generalized Bloch theorem the rotating-frame
Bloch Hamiltonian is
\begin{equation}
  H(k)=\begin{pmatrix}\varepsilon(k-\tfrac{q}{2}) & -\Delta\\ -\Delta & \varepsilon(k+\tfrac{q}{2})\end{pmatrix},
  \qquad \varepsilon(k)=-2t\cos(kd),
\end{equation}
matching the paper's up/down plane-wave split (their Eq.~7). The charge velocity
is $\partial_k H$, the spin-flux tensor is $Q=\mathrm{Re}\langle\psi^\dagger
(S\otimes v)\psi\rangle$ (their Eq.~1), and the torque--current tensor $C$ is
built from the constant-$\tau$ FS integral (Eqs.~3--6; $\tau$ cancels).

\section{Machine-checkable claims and results}
\begin{table}[h]\centering
\begin{tabular}{clllc}
\toprule
\# & Claim & Paper & Computed & Verdict\\
\midrule
C1 & $C$ structure (axis current $\to$ single & single nonzero & rotate $=0.0282$; & \pass\\
   & rotate-spiral torque; out-of-plane $=0$) & ($C_{23}$ only) & out-of-plane $\sim\!10^{-16}$; $S_z\!\approx\!0$ & \\
C3 & $|P|=0.5\Rightarrow$ tilt from axis & \ang{30} & \ang{30.000} & \pass\\
C4 & $q=0.20\,(2\pi/c)\Rightarrow$ phase/layer & $0.20\pi=0.6283$ & $0.6283$\,rad & \pass\\
C5 & rotation freq $\propto j$ & linear & dev.\ $3.5\times10^{-18}$ & \pass\\
C6 & crude/microscopic ratio & $\sim4\times$ & $0.155$ (model-dep.) & \partialv\\
-- & abs.\ freq at \SI{e7}{A/cm^2} & \SI{0.07}{GHz} & not recomputed (DFT) & \oos\\
\bottomrule
\end{tabular}
\end{table}

\noindent\textbf{Machine-checkable pass rate: 4/5.}

\subsection*{C1 --- tensor structure (PASS)}
For the planar spiral the axis current drives a single dominant in-plane
spin-flux channel while the out-of-plane ($z$) spin flux vanishes to machine
precision. The two in-plane channels came out numerically \emph{equal}
($S_x=S_y=0.0282$) with $S_z\approx-2\times10^{-16}$, a clean planar-spiral
signature that reproduces the paper's single-nonzero-component $C$: only a
current along the axis rotates the spiral.

\subsection*{C3/C4 --- geometry (PASS)}
$|P|=0.5\Rightarrow\arcsin(0.5)=\ang{30}$ tilt; $q=0.20\,(2\pi/c)$ with atomic
layer spacing $c/2$ gives a per-layer transverse-spin phase advance of
$0.20\pi$\,rad, matching the paper's ``$q\pi$ per layer.''

\subsection*{C5 --- linear scaling (PASS)}
The induced rotate-spiral torque is strictly proportional to $j$ over
$10^{9}$--$10^{12}\,\si{A/m^2}$ (deviation $\sim10^{-18}$), consistent with the
bulk linear-response STT picture (above the anisotropy critical current).

\subsection*{C6 --- crude vs.\ microscopic ratio (PARTIAL, honest negative)}
The paper's crude adiabatic estimate is $\sim4\times$ its microscopic $C$-matrix
value (``catches the order of the effect''). In our toy model the crude estimate
\emph{under}-estimates (ratio $0.155$), not $4\times$. Root cause: the crude
prefactor and the microscopic $Q$-tensor use different FS weightings and band
counts; the paper's factor-4 is order-of-magnitude for Er's specific multi-band
nesting FS and is not expected to survive in a single-band 1D model. Reported
as-is, not tuned (see \texttt{failure\_analysis.md}).

\section{Verdict}
\textbf{The paper's central, convention-independent claims replicate.} A charge
current along a planar spin-spiral axis produces a single-component
torque--current tensor that rigidly rotates the spiral, with the correct
tilt/per-layer geometry and strict linear current scaling. The material-specific
DFT numbers (\SI{0.07}{GHz}, $C_{23}=0.5\,\hbar\,\si{\angstrom^2}$) are out of
scope for this offline replication and were not recomputed. One internal
consistency check (the $\sim4\times$ crude/microscopic ratio) is model-dependent
and did not reproduce numerically --- reported honestly.

\medskip
\noindent\textbf{Coverage: 7/10.} All qualitative/structural/scaling claims and
the geometry are covered by real runnable code; the DFT-specific absolute
magnitude and the numeric $4\times$ factor are not.

\noindent\textbf{Agreement: 8/10.} Every covered claim agrees quantitatively
(exact for geometry/structure/linearity); the single miss (C6 numeric factor) is
understood and expected for a toy model, not a contradiction of the paper.

\end{document}
